\documentclass[prc,amsart,english]{revtex4}
\usepackage{graphicx}
\usepackage{epsfig}
\usepackage{bm}
\usepackage{color}
\usepackage{float}
\usepackage{dcolumn}
\usepackage{multirow} 
\usepackage{hyperref}

\usepackage[T1]{fontenc}       % DC-fonts
\begin{document}
\title{Nuclear Talent Course on Nuclear Theory for Nuclear Structure Experiments, ECT*, July 3-21, 2017}
\maketitle

\includegraphics[width=0.9\columnwidth]{fig1.jpg}

\clearpage

The present document aims at giving a summary, 
with background information as well, of  the twelfth Nuclear Talent course, {\em Nuclear Theory for Nuclear Structure Experiments}.
The course was held at the premises of the European Center for Theoretical Studies
in Nuclear Physics and Related Areas (ECT*), Trento, Italy during the period July 3 to July 21, 2017. 

This was the fourth course which has been held at the ECT*, the first one in 2012, then 2014 and 2015 and finally this one in 2017. The support from the ECT* has simply been crucial for the successful arrangement and running of these courses. In total 52 applications were received and a final of 31 participants were selected. Five teachers were involved, with two of them being the organizers. 


\section{Report }

\subsection{Introduction}
The nuclear shell model plays  a central role in guiding our analysis of the wealth of nuclear structure 
experimental data.
It provides an excellent link to the underlying nuclear forces and the pertinent laws of motion,
allowing nuclear physicists to interpret complicated experiments in terms of various 
components of the nuclear Hamiltonian and to understand a swath of nuclei by following chains of isotopes and isotones over wide ranges of nucleon numbers. The nuclear shell model allows us to see how the structure of nuclei changes and how the occupation of specific nucleonic orbits affects the interplay of residual interactions and configuration mixing.  The computed expectation values and transition likelihoods 
can be directly linked to experiment, with the potential to single out new phenomena and guide future experiments. 
Large-scale shell-model calculations represent also 
challenging computational and theoretical topics, spanning from efficient usage of high-performance computing facilities  to consistent theories for deriving effective Hamiltonians and operators. 
Alltogether, these various facets of nuclear theory represent important elements in our endeavors to understand nuclei and their limits of stability.  
The goal and motivation of this course was to introduce and
develop the nuclear structure tools  needed to carry out forefront research using the shell model as the central tool. The various exercises and project focussed on the development of a shell-model code for simpler systems like $sd$-shell nuclei, giving the participants the essential ideas of configuration interaction methods. During the first two weeks the aim was to develop such a shell-model code. With these insights, the students could divert into several directions
the last week, from the usage of the NushellX suite of nuclear structure programs to further developing 
their own shell-model program. After completion, it is our sytong feeling that the participants have understood the overarching ideas 
behind central theoretical tools used to analyse nuclear structure experiments.


\subsection{Aims and outcomes}

We held approximately forty-five hours of lectures over three weeks and a
comparable amount of practical computer and exercise sessions, including the
setting of individual problems and the organization of exercises and projects.

The mornings consisted of lectures and the
afternoons was  devoted to exercises and project work meant to shed light on the exposed theory, 
the computational project  and eventual individual student projects.

The total load is approximately 170 hours, corresponding to  {\bf 7 ECTS} in Europe.
The days were organized as follows: 9.30-13 lectures, time for exercises with assistance (including lunch)
till 18-19 (3-4 hours of allocated exercise sessions per day).
The course was held at the premises ECT* (Villazzano) in Trento, Italy, from
July 3  to July 21 in 2017.
Professor Takaharu Otsuka gave a special lecture on Shell-Model Monte Carlo methods in the second week and 
in addition, five students gave presentations of their own research. These lectures were closely related to the topics of the present Talent course.  



All learning material, with preparatory material on second quantization plus Alex Brown's lecture notes were provided well in advance before the school (three weeks) and available via the Github link \href{https://github.com/NuclearTalent/NuclearStructure}. 

The students were divided into nine groups from the very first day of the school. These groups remained stable throughout the whole duration of the course. Their basic set up 
consisted in having at least one experimentalist and one theorist in every group. Three groups had four members  while the remaining six groups consisted of three members only. 
This can be summarized in
\begin{itemize}
\item We  used group work to carry out the tasks.

\item Groups where carefully put together to maximize diversity of backgrounds.  

\item The weekly results and progress of every group was presented at the end of week (Friday afternoon) in a conference-like setting to create accountability.  

\item We several organized events (weekly social dinners as well) where individuals and groups exchange their experiences, difficulties and successes to foster interaction.  

\item During the school, on-line and lecture-based training tailored to technical issues were provided. The participants learnt to use and interpret the results of computer-based and hand calculations of nuclear models. The lectures were aligned with the practical computational projects and exercises and the lecturers were always available to help students and work with them during the project and exercise sessions.  

\item During the second week, after having introduced the basic theory needed to understand full configuration interaction  theory (the basis for shell-model studies), the theoretical interpretations were accompanied by talks on how to perform nuclear structure experiments. These lectures were given by Alexandra Gade and Robert Grzywacz and illustrated nicely how nuclear structure theory can be used  to interpret data. Several of the experimental students gave additional lectures the same week on additional experimental techniques. Overall, these lectures illustrated nicely, and were highly appreciated by the students, how tightly interwoven theory and experiment are.
  
\item There was enough flexibility and time to organize mini-lectures and discussion sessions on an ad-hoc basis.  

\item Each group of students maintained an online logbook of their activities and results using Github. This allowed the teaching team to monitor the weekly progress of each group. The weekly presentations were also stored at the Github address of each group. Each group had to establish a Github address where all material was placed.  

\item Training modules, codes, lectures, practical exercise instructions, online logbooks, instructions and information created by participants are stored  into a comprehensive website that is available to the community and the public for self-guided training or for use in various educational settings (for example, a graduate course at a university could assign some of the projects as homework or an extra credit project, etc). The website with all material is at \href{https://github.com/NuclearTalent/NuclearStructure}. 

\item Seven students handed in their final reports and requested credits. 
  
\end{itemize}


\subsection{Aims and Learning Outcomes}
At the end of the course, our ideal was that the students should have a basic understanding of
\begin{itemize}
\item Configuration interaction methods (nuclear shell-model here) as a central tool to interpret nuclear structure experiment

\item Have an understanding of single-particle basis functions and the construction of many-body basis states built thereupon. Examples are basis states from a Woods-Saxon potential, harmonic oscillator states and mean-field based states from a Hartree-Fock calculation. The single-particle basis states are orthonormal and are used to construct  a corresponding orthonormal basis set of Slater determinants.

\item Develop an understanding of what defines an observable. 

\item Understand how theory  can be used to interpret experimental quantities (separation energies and shell gaps for example).

\item Understand how experiment is used to extract transition probabilities and information about ground and exited states. 

\item Understand how second-quantization is used to represent states and compute expectation values and transition probabilities  of operators

\item Understand how the Hamiltonian matrix is constructed from this orthonormal basis set of many-body states (linear expansion of Slater determinants)

\item For nuclear systems like the $sd$-shell, essentially all nuclei can be studied using direct diagonalization methods. In the construction of the shell-model code during the first two weeks, the students should learn to compute  the ground state and the excited states of selected $sd$-shell nuclei. This project applies to all students. During the last week, students can pursue more individually defined projects

\item The students will also learn to understand the basic elements of effective shell-model Hamiltonians and how to interpret the calculated properties in terms of various components of the nuclear forces (spin-orbit force, tensor force, central force etc). We provided the students with the necessary tools to perform such analyses. 

\item Understand how to use shell-model calculations to calculate decay rates and transition probabilities and relate these to various electromagnetic transition operators and operators for beta-decays and double-beta decays.  

\item Develop a critical understanding the limits of shell-model studies and how these can be related to interpretations of data such as results from in-beam and decay experiments.

\item Understand how to use second-quantization to construct one-body and two-body transfer operators, overlap functions, spectroscopic factors and experiments related to spectroscopic factors. 

\item For the students which wish to follow a more computational path during the last week, iterative eigenvalue solvers will discussed. Similarly, efficient representations of many-body states and computations of Hamiltonian matrices were also be discussed.

\item We will also discuss modern shell-model codes like NushellX. This suite of programs can be used by students to pursue their own projects. Applications of NushellX to the calculations of various observables were discussed.
\end{itemize}

Based on the feedback from the students, see the end of the report, we feel these learning outcomes were achieved to a large extent. All groups developed their own shell-model code and benchmarked it with results from NushellX for $sd$-shell nuclei. Furthermore, NushellX was used to perform advanced interpretations of experimental data, allowing thereby the participants to both understand the basic theory (and later if of interest develop their own advanced nuclear structure tools) and to relate theory to experimental data from various decay modes.

All in all, for both teachers and students, we believe most of our goals were achieved. 

\subsection{Course Schedule}
  
Lectures were approximately 45 min each with a small break between each lecture. 
The acronyms are 
\begin{itemize}
\item AG: Alexandra Gade

\item BAB: Alex Brown

\item GJ: Gustav Jansen

\item MHJ: Morten Hjorth-Jensen

\item RG: Robert Grzywacz
\end{itemize}
The final schedule is listed here. 
\noindent
\paragraph{Week 1, July 3-7.}
\begin{footnotesize}
\begin{quote}
\begin{tabular}{llll}
\hline
\multicolumn{1}{c}{ Day } & \multicolumn{1}{c}{  } & \multicolumn{1}{c}{ Lecture Topics and lecturer } & \multicolumn{1}{c}{ Projects and exercises } \\
\hline
Monday 3    & 9am-930am     & Registration at the ECT                                             &                                                                  \\
            & 930am-10am    & Introduction and welcome (BAB, MHJ and GJ)                          &                                                                  \\
            & 10am-1230pm   & Survey of data (BAB)                                                & Intro to NushellX (BAB)                                          \\
            & 1230pm-230pm  & Lunch +own activities                                               &                                                                  \\
            & 230pm-6pm     & More survey of data (BAB)                                           & NushellX hands-on and discussion of project (MHJ)                \\
\hline
Tuesday 4   & 930am-1130pm  & Mean-field and shell-model basics (MHJ)                             & Discussion of simple shell-model project \\
            & 1130am-1230pm & Single-particle potentials and shell-model basics (BAB)             &                                                                  \\
            & 1230pm-230pm  & Lunch+ own activities                                               &                                                                  \\
            & 230pm-6pm     & Single-particle potentials                                          & Studies of separation energies, shell gaps etc                   \\
\hline
Wednesday 5 & 930am-1130am  & Welcome by ECT director Jochen Wambach and Shell-model basics (MHJ) & Discussion of simple shell-model project                         \\
            & 1130am-1230pm & Shell-model dimensionalities (BAB)                                  &                                                                  \\
            & 1230pm-230pm  & Lunch+ own activities                                               &                                                                  \\
            & 230pm-6pm     &                                                                     & Work on shell-model project and NushellX                         \\
\hline
Thursday 6  & 930am-1130am  & Shell-model basics (MHJ)                                            & Discussion of simple shell-model project                         \\
            & 1130am-1230pm & Proton-neutron formalism and isospin (BAB)                          &                                                                  \\
            & 1230pm-230pm  & Lunch+ own activities                                               &                                                                  \\
            & 230pm-6pm     &                                                                     & Work on shell-model project and NushellX                         \\
\hline
Friday 7    & 930am-1230pm  & Effective interactions for shell model (GJ)                         &                                                                  \\
            & 1230pm-230pm  & Lunch+ own activities                                               &                                                                  \\
            & 230pm-5pm     &                                                                     & Group presentations   \\
\hline
\end{tabular}
\end{quote}
\end{footnotesize}
\noindent

\paragraph{Week 2, July 10-14.}
\begin{footnotesize}
\begin{quote}
\begin{tabular}{llll}
\hline
\multicolumn{1}{c}{ Day } & \multicolumn{1}{c}{  } & \multicolumn{1}{c}{ Lecture Topics and lecturer } & \multicolumn{1}{c}{ Projects and exercises } \\
\hline
Monday 10    & 930am-1030am  & Masses from experiment  (AG)                                                                          &                                                                \\
             & 1030am-1230pm & Effective Hamiltonians (BAB and GJ)                                                                   &                                                                \\
             & 1230pm-230pm  & Lunch +own activities                                                                                 &                                                                \\
             & 230pm-6pm     &                                                                                                       & Work on shell-model project and NushellX                       \\
\hline
Tuesday 11   & 930am-1030am  & Spectroscopic factors from experiment (AG)                                                            &                                                                \\
             & 1030am-1230pm & Spectroscopic factors from theory (BAB)                                                               &                                                                \\
             & 1230pm-230pm  & Lunch+ own activities                                                                                 &                                                                \\
             & 230pm-330pm   & Monte Carlo Shell-model seminar by Takaharu Otsuka                                                    &                                                                \\
             & 330pm-6pm     &                                                                                                       & Work on shell-model project and NushellX                       \\
\hline
Wednesday 12 & 930am-1130am  & Introduction to $\beta$-decay, experiment (RG)                                                        &                                                                \\
             & 1130am-1230pm & $\beta$-decay and sum rules from theory (BAB)                                                         &                                                                \\
             & 1230pm-230pm  & Lunch+ own activities                                                                                 &                                                                \\
             & 230pm-330pm   & Discussion of IBM and shell-model by Noam, Tobias and Takaharu                                        &                                                                \\
             & 330pm-6pm     &                                                                                                       & Work on shell-model project and NushellX                       \\
\hline
Thursday 13  & 930am-1030am  & Electromagnetic decays from experiment (AG)                                                           &                                                                \\
             & 1030am-1230pm & Shell-model transition probabilities (BAB)                                                            &                                                                \\
             & 1230pm-230pm  & Lunch+ own activities                                                                                 &                                                                \\
             & 230pm-330pm   & Presentation by Guillaume                                                                             &                                                                \\
             & 330pm-6pm     &                                                                                                       & Work on shell-model project and NushellX                       \\
\hline
Friday 14    & 930am-11am    & Advanced aspects of  $\beta$-decay, experiment (RG)                                                   &                                                                \\
             & 11am-1230pm   & Summary of first two weeks  (MHJ) &                                                                \\
             & 1230pm-230pm  & Lunch+ own activities                                                                                 &                                                                \\
             & 230pm-5pm     &                                                                                                       & Group presentations \\
\hline
\end{tabular}
\end{quote}
\end{footnotesize}
\noindent

\paragraph{Week 3, July 17-21.}

\begin{footnotesize}
\begin{quote}
\begin{tabular}{llll}
\hline
\multicolumn{1}{c}{ Day } & \multicolumn{1}{c}{  } & \multicolumn{1}{c}{ Lecture Topics and lecturer } & \multicolumn{1}{c}{ Projects and exercises } \\
\hline
Monday 17    & 930am-12pm   & Building a shell-model code (MHJ)                                                                              &                                                                \\
             & 12pm-1pm     & More on $\beta$-decay and $\gamma$ decay  (BAB)                                                                &                                                                \\
             & 1pm-230pm    & Lunch +own activities                                                                                          &                                                                \\
             & 230pm-6pm    &                                                                                                                & Work on shell-model project and NushellX                       \\
\hline
Tuesday 18   & 930am-12pm   & Building a shell-model code, Lanczos' algorithm and Hartree-Fock theory (MHJ)                                  &                                                                \\
             & 12pm-1pm     & $\gamma$ decay  (BAB)                                                                                          &                                                                \\
             & 1pm-230pm    & Lunch+ own activities                                                                                          &                                                                \\
             & 230pm-330pm  & Laser spectroscopy talk by Agi and Shane                                                                       &                                                                \\
             & 330pm-6pm    &                                                                                                                & Work on shell-model project and NushellX                       \\
\hline
Wednesday 19 & 930am-12pm   & From Hartree-Fock theory to the shell-model, analyzing the results (MHJ)                                       &                                                                \\
             & 12pm-1pm     & $\gamma$ decay Two-body transition operators and double-beta decay (BAB)                                       &                                                                \\
             & 1pm-230pm    & Lunch+ own activities                                                                                          &                                                                \\
             & 230pm-330pm  & No-core shell-model talk by Patrick  and Nuclear interactions by Tor                                           &                                                                \\
             & 330pm-6pm    &                                                                                                                & Work on shell-model project and NushellX                       \\
\hline
Thursday 20  & 930am-1230pm & From Hartree-Fock theory to the shell model, discussion of final project  and summmary of course (BAB and MHJ) &                                                                \\
             & 1230pm-230pm & Lunch+ own activities                                                                                          &                                                                \\
             & 230pm-5pm    &                                                                                                                & Group presentations \\
\hline
Friday 21    & 930am1230pm  & Project work and own activities (no lectures)                                                                  &                                                                \\
             & 1230pm-230pm & Lunch+ own activities                                                                                          &                                                                \\
             & 230pm-6pm    & Own activities                                                                                                 &                                                                \\
\hline
\end{tabular}
\end{quote}
\end{footnotesize}



\subsection{Teachers and organizers}


The organizers were

\begin{enumerate}
\item \href{{https://people.nscl.msu.edu/~brown/}}{Alex Brown}, \href{{http://www.nscl.msu.edu/}}{National Superconducting Cyclotron Laboratory} and \href{{https://www.pa.msu.edu/}}{Department of Physics and Astronomy}, \href{{http://www.msu.edu/}}{Michigan State University}, East Lansing, MI 48824, USA

\item \href{{http://mhjgit.github.io/info/doc/web/}}{Morten Hjorth-Jensen}, \href{{http://www.nscl.msu.edu/}}{National Superconducting Cyclotron Laboratory} and \href{{https://www.pa.msu.edu/}}{Department of Physics and Astronomy}, \href{{http://www.msu.edu/}}{Michigan State University}, East Lansing, MI 48824, USA and Department of Physics, University of Oslo, N-0316 Oslo, Norway 
\end{enumerate}

\noindent
Morten Hjorth-Jensen will also function as student advisor and coordinator.

The teachers were

\begin{enumerate}
\item \href{{https://people.nscl.msu.edu/~brown/}}{Alex Brown}, \href{{http://www.nscl.msu.edu/}}{National Superconducting Cyclotron Laboratory} and \href{{https://www.pa.msu.edu/}}{Department of Physics and Astronomy}, \href{{http://www.msu.edu/}}{Michigan State University}, East Lansing, MI 48824, USA

\item \href{{https://people.nscl.msu.edu/~gade/}}{Alexandra Gade}  \href{{http://www.nscl.msu.edu/}}{National Superconducting Cyclotron Laboratory} and \href{{https://www.pa.msu.edu/}}{Department of Physics and Astronomy}, \href{{http://www.msu.edu/}}{Michigan State University}, East Lansing, MI 48824, USA

\item \href{{http://web.utk.edu/~rgrzywac/}}{Robert Grzywacz}  at \href{{http://www.ornl.gov/}}{Oak Ridge National Laboratory}, Oak Ridge, TN 37831  and \href{{https://www.phys.utk.edu/}}{Department of Physics and Astronomy}, \href{{http://www.utk.edu/}}{University of Tennessee}, Knoxville, TN 37996-1200, USA

\item \href{{http://mhjgit.github.io/info/doc/web/}}{Morten Hjorth-Jensen}  at \href{{http://www.nscl.msu.edu/}}{National Superconducting Cyclotron Laboratory} and \href{{https://www.pa.msu.edu/}}{Department of Physics and Astronomy}, \href{{http://www.msu.edu/}}{Michigan State University}, East Lansing, MI 48824, USA and  Department of Physics, University of Oslo, N-0316 Oslo, Norway

\item \href{{https://www.ornl.gov/staff-profile/gustav-r-jansen}}{Gustav Jansen}  at \href{{http://www.ornl.gov/}}{Oak Ridge National Laboratory}, Oak Ridge, TN 37831, USA

\end{enumerate}


\subsection{Budget}
The teachers were mostly self-sponsored, with Gade, Brown, Grzywacz and Jansen being supported by the FRIB theory alliance, Oak Ridge National Laboratory and the University of Oslo. The housing of Gade and Jansen was supported by the ECT* since they were participants of a workshop which was held during the second week of the Talent course. The ECT* covered lunches for all teachers during weekdays. Hjorth-Jensen was fully self-supported (except for lunches from the ECT*). Furthermore, Jansen's travel was piad by the University of Oslo. 

The ECT* covered the local expenses (housing and food) for 19 students, totaling 19000 Euros. Of the 31 students, three students came from the University of Trento or were residents of Trento and received only lunches from the ECT*. The remaining nine students were self-supported, but lunches during the weekdays were provided by the ECT*. 

The ECT* provided support of a total of 19000 Euros, plus lunches for 15 days for 31 students and five teachers. 
The estimated price per lunch is 10-15 Euros, totaling thus approximately 6000 Euros.
The total direct contributions from the ECT*, including housing for Jansen and Gade as well as Otsuka,  amount thus approximately to at least 30000 Euros. This excludes usage of the premises of the ECT*, the fantastic administrative support from the ECT* (Serena is just amazing) and other running expenses. Here we would estimate these to amount to approximately  10000 Euros. We find it reasonable to assume that the total support received from the ECT*, direct and indirect expenses, amount to roughly 40000 Euros. 


All travels were covered by the home institutions of the participants. If we assume on average an estimate of 1000 Euros per travel per participants, and that the teachers were essentially fully self-supported, plus nine fully self-supported expenses, it is fair to assume that the additional expenses amount to approximately $36\times 1000=36000$ Euros for travel, $9\times 1000=9000$ Euros for the self-supported students and approximately 10000 Euros to cover housing and additional living expenses for the teachers. In total, a realistic estimate of the
budget is close to 90000-100000 Euros, with the ECT* sponsoring at least 40000 Euros. 

This estimate reflects  (without an honorarium for the teachers) a realistic expenditure for a three week intensive course with 31 students.  Housing expenses will however vary from place to place. The ECT* has a long tradition in running doctoral schools and a good practice on finding reasonable housing for the participants. If administrative help is to be added to the budget, one needs at least to account for 
three full weeks of administrative help from one person.



\subsection{Participants with their home institutions and nationalities}
The target group for the Nuclear Talent courses is Master of Science students, PhD students and early post-doctoral fellows.
Out of 52 applicants, 31 students were selected. Priority was given
to Master of Science students and early PhD students. 

The students were expected to have operating programming skills in Fortran/C++/Python and knowledge of
quantum mechanics at an intermediate level, with basic knowledge of 
many-body physics. 
Of the 31 students, eight were Master of Science students and all, if possible,  expressed they would like to continue with a PhD in nuclear physics. Of the PhD students, the majority
were in their first two years. Sixteen nationalities (and three continents) and 20 different institutions/affiliations 
were represented. Twelve out of the 31  students were experimentalists and nine were female students. 
The students and their respective institutions are listed in the table below. 
\begin{table}[hbtp]
  \begin{ruledtabular}
    \begin{tabular}{|c||l|l|l|l|}
     Name & Level &Institution &Nationality & Theory/Expt  \\\hline
     Ashley Hood(F)& PhD &Lousiana State University, USA  &USA & Expt\\
     Erin Good(F)& PhD & Lousiana State University, USA &USA & Expt \\
     Anne Marie Forney(F)&PhD  &University of Maryland, USA  &USA & Expt\\
     Patrick Fasano(M)& PhD & Notre Dame University, USA &USA & Theory \\
     Guillaime H\"{a}fner(M)& MSc & University of Cologne, Germany &Germany & Expt \\
     Udo Gayer(M)& PhD & TU Darmstadt, Germany & Germany & Expt \\
     Rosa-Belle Gerst(F)& PhD & University of Cologne, Germany & Germany &Expt\\
     Tobias Beck(M)& PhD & TU Darmstadt, Germany &Germany &Expt\\
     Paolo Andreatta(M)& PhD & University of Trento, Italy  &Italy& Theory \\
     Luca Riz(M)& PhD & University of Trento, Italy &Italy & Theory \\
     Giovanni Pederiva(M)& MSc & University of Oslo, Norway  &Italy &Theory \\
     Gianluca Salvioni(M)& PhD &University of Jyv\"askyl\"a, Finland  &Italy& Theory   \\
     David Miur(M)& PhD & University of York, UK  &UK & Theory\\
     Matthew Shelley(M)& MSc &University of York, UK &UK & Theory \\
     Shane Wilkins(M)& PhD & University of Manchester, UK  &UK & Expt \\
     Ina Berentsen(F)& MSc & University of Oslo, Norway &Norway & Expt \\
     Mathias Vege(M)& MSc & University of Oslo, Norway &Norway& Theory \\
     Sean Miller(M)&MSc  &University of Oslo, Norway  &Norway& Theory \\
     Tor Dj\"{a}rv(M)&PhD  &Chalmers  &Sweden & Theory \\
     Martin Albertsson(M) &PhD  &Lund University, Sweden  &Sweden & Theory \\
     Lucia Campo Crespo(F)& PhD & University of Oslo, Norway  &Spain & Expt\\
     Antonio Marquez Romero(M)&PhD  &University of York, UK  &Spain & Theory \\
     Gilho Ahn(M) & MSc & University of Athens, Greece & Greece & Theory \\
     Ovidiu Nitescu(M) & MSc & University of Bucuresti, Romania & Romania & Theory \\
     Noam Gavrielov(M) & PhD & The Hebrew University, Jerusalem, Israel & Israel & Theory \\
     Olvier Vasseur(M) & PhD & IPN Orsay, Paris, France & France & Theory \\
     Liqiang Qi(M) & PhD & IPN Orsay, Paris, France & China & Expt \\
     Mahsa Razazi(F) & PhD & Islamic Azad University, Urmia, Iran & Iran & Theory \\
     Masaaki Tokiada(M) & PhD & University of, Japan & Japan & Theory \\
     Agota Kozsorus(F) & PhD & University of Leuven, Belgium & Hungary & Expt \\
     Tiia Haaverinen(F)& PhD &University of Jyv\"askyl\"a, Finland  &Finland & Theory   \\ \hline
    \end{tabular}
  \end{ruledtabular}
\end{table} 

On average, the students performed very well. It was an excellent  and very active group. Seven students handed in the final report and will receive a credit transfer of the seven ECTS from the University of Trento.

\clearpage

\section{Brief summary}

The support from the ECT*, and its year-long experiences with running doctoral training programs and Talent courses, was central to the success of the course. Of uttermost importance was Serena Degli Avancini's help will all administrative matters, from housing to minor practicalities.  Without her help and the other staff of the ECT*, and the director's (Prof.~Jochen Wambach) 
enthusiastic support,  it is unlikely that this course would have run as smoothly as it did. 

As a recommendation  to new teaching teams, setting student groups from day one is essential. Groups of 3-4 students were optimal and students with complementary and/or different skills interest did benefit from each other. This group of students was particularly well functioning and setting up working gropus from the very beginning helped in developing new friendships and perhaps even new collaborations. The feedback from the students really testifies to this.

Else, having proper projects and exercises which are coordinated with the lectures is essential. The learning outcomes and the feeling of having written one's own shell-model code cannot be underestimated.

With several  teachers it is also difficult to achieve a uniform teaching style. And some of the lectures were best suited in terms of slides. Other teachers used a mix of slides and blackboard, or just the blackboard. The latter slows down the pace and many of the respondents seemed to prefer that teaching modus. But with several teachers, it is also important that each teacher feels comfortable with her/his way of teaching, especially since this was an intensive three weeks course.

Overall, we as teaching feel the original learning outcomes and goals were met. It was a marvelous group of students to teach. 
\noindent\bigskip
\includegraphics[width=0.9\columnwidth]{fig2.jpg}
\end{document}

\documentclass[superscriptaddress,amsmath,amssymb,aps,floatfix]{revtex4-2}
\usepackage{graphicx,tikz} 
\usepackage[mode=buildmissing]{standalone} 
\usepackage{dcolumn}
\usepackage{bm}
\usepackage{physics}
\usepackage{siunitx}
\usepackage{soul}
\usepackage{comment}
\usepackage[bb=boondox]{mathalfa}
\usetikzlibrary{external}
\tikzexternalize[prefix=figures/]
\DeclareMathOperator{\sinc}{sinc}
\newcommand{\may}[1]{\textcolor{orange}{#1}}
\newcommand{\ale}[1]{{\color{blue}{\bf Ale}: #1}}
\newcommand{\zd}[1]{{\color{purple}{\bf ZD}: #1}}
\usepackage{xcolor}

\begin{document}
\title{Proposal for a Nuclear Talent course at the ECT* in 2025:\\Quantum Computing for Nuclear
Physics}


\author{Alexei Bazavov}

\affiliation{Department of Computational Mathematics, Science and
Engineering and Department of Physics and Astronomy, Michigan State
University, East Lansing, MI 48824, USA}

\author{Zohreh Davoudi}
\affiliation{Department of Physics, University of Maryland, College Park,
MD 20742, USA}
\affiliation{Center for Quantum Information and Computer Science (QuICS), University of Maryland, College Park, MD 20742, USA}

\author{Morten Hjorth-Jensen}
\affiliation{Department of Physics and Center for Computing in Science Education, University of Oslo, N-0316 Oslo, Norway}
\affiliation{Facility for Rare Isotope Beams and Department of Physics and Astronomy, Michigan State University, East Lansing, MI 48824, USA}

\author{Ryan LaRose}
\affiliation{Department of
Computational Mathematics, Science and Engineering, Michigan State
University, East Lansing, MI 48824, USA}

\author{Dean Lee}
\affiliation{Facility for Rare Isotope Beams and Department of Physics and
Astronomy, Michigan State University, East Lansing, MI 48824, USA}

\author{Alessandro Roggero} 
\affiliation{Department of Physics, University of Trento, Povo, 38123 Trento, Italy}
    
\maketitle

\section{Proposal for a Nuclear Talent course at the ECT* in 2025}

We would like to propose a three-week Nuclear Talent course on theory
for exploring quantum computing and quantum information science
applied to nuclear physics at the ECT* for the summer of 2025. We
believe such a course has a strong potential to attract many students,
theorists and experimentalists alike.

Below we give a short motivation for the proposed course and the
rationale behind the Nuclear Talent initiative. Thereafter we detail
our course plans with learning outcomes, objectives and teaching
philosophy, as well as various organizational and practical matters.

The teaching teams consists of quantum computing theorists and nuclear
theorists, with expertise ranging from many-body methods to quantum
field theories. Several of us have developed courses on quantum
computing and/or taught similar courses, in addition to our ongoing
research on quantum computing. This spans from quantum engineering to
developments of new algorithms and error correction studies. We
believe such a mix is important as it gives the students a better
understanding on how quantum computing can be applied to nuclear
physics problems, and what are the limitations and possibilities in
understanding and interpreting the various algorithms on present and
planned quantum technologies.

\subsection{Motivation}\label{motivation}

For nuclear theorists, the overarching challenge is to develop a
comprehensive description of nuclei and their reactions, grounded in
the fundamental interactions between the constituent nucleons with
quantifiable uncertainties to maximize predictive power. As
experimental frontiers have shifted to the study of rare isotopes, the
predictive power of successful phenomenological approaches like the
shell model and density functional theory is challenged by the
scarcity of nearby experimental data to constrain model
parameters. Therefore, it is expected that few-body and many-body
methods will play an increasingly prominent role to help improve the
predictive power of such ``data driven'' methods as experiment moves
deeper into largely unexplored regions of the nuclear chart.

To understand why nuclear matter is stable, and thereby shed light on
the limits of nuclear stability, is one of the overarching aims and
intellectual challenges of basic research in nuclear physics. To
relate the stability of matter to the underlying fundamental forces
and particles of nature as manifested in nuclear matter, is central to
present and planned rare isotope facilities. From a theoretical
standpoint, this involves understanding how the basic building blocks
of Nature interact and conspire to build up atomic nuclei as we know
them, with the aim to understand what makes visible matter stable. The
theoretical efforts span from methods like lattice quantum
chromodynamics, via effective field theories to many-body theories
applied to atomic nuclei and infinite nuclear matter. All these
methods rely on theoretical approximations whose applicabilities are
often limited by the dimensionality of the specific problem being
studied. In recent years, there has been considerable progress in
developing quantum-computing algorithms applied to quantum many-body
systems, with the hope to circumvent many of the classically
intractable problems.

This proposal for a Nuclear Talent school aims at bringing together
the efforts of nuclear many-body theorists, quantum information
theorists, and mathematicians in order to present and discuss
algorithms for studying nuclear systems using recent progress in
quantum information theory.

\subsection{Introduction to the Talent Courses}

The TALENT initiative, \href{http://www.nucleartalent.org}{Training in
  Advanced Low Energy Nuclear Theory}, aims at providing an advanced
and comprehensive training to graduate students and young researchers
in low-energy nuclear theory. The initiative is a multinational
network of several European and Northern American institutions and
aims at developing a broad curriculum that will provide the necessary
training in cutting-edge theory for understanding nuclei and nuclear
reactions.  These objectives will be met by offering a series of
lectures, delivered by experts in nuclear many-body theory and quantum
information theories.  The educational material generated under this
program will be collected in the form of WEB-based courses, textbooks,
and a variety of modern educational resources. No such
all-encompassing material is available at present; its development
will allow dispersed university groups to profit from the best
expertise available worldwide.

The Nuclear Talent initiative has organized and run several advanced
courses since the summer of 2012. Several of these courses have been run
and organized (in a very successful way) at the premises of the
ECT\emph{. We hope thus, if this course gets approved by the board of
directors of the ECT}, that we can continue this successful and very
fruitful collaboration.

\section{Aims and Learning Outcomes}
\subsection{Format}

We propose approximately forty-five hours of lectures over three weeks
and a comparable amount of practical computer and exercise sessions with
supervised practices and tutorials.

The mornings will consist of lectures and the afternoons will be devoted
to exercises meant to shed light on the exposed theory, the
computational projects, and individual student projects. These components
will be coordinated to foster student engagement, maximize learning, and
create lasting value for the students. For the benefit of the TALENT
series and of the community, material (courses, slides, problems and
solutions, reports on students' projects) will be made publicly
available using version control software like \emph{git} and posted
electronically on \href{https://github.com}{github}.

As with previous TALENT courses, we envision the following features for
the afternoon sessions: We will use both individual and group work to
carry out tasks that are very specific in technical instructions, but
leave freedom for creativity.

\begin{itemize}
\item
  Groups will be carefully put together to maximize diversity of
  backgrounds.
\item
  Results will be presented in a conference-like setting to create
  accountability.
\item
  We will organize events where individuals and groups exchange their
  experiences, difficulties, and successes to foster interaction.
\item
  During the school, on-line and lecture-based training tailored to
  technical issues will be provided. Students will learn to use and
  interpret the results of computer-based and hands-on calculations of
  quantum computing algorithms. The lectures will be aligned with the
  practical computational projects and exercises and the lecturers will
  be available to help students and work with them during the exercise
  sessions.
\item
  These interactions will raise topics not originally envisioned for the
  course but which are recognized to be valuable for the students. There
  will be flexibility to organize mini-lectures and discussion sessions
  on an ad-hoc basis in such cases.

\item
  Training modules, codes, lectures, practical exercise instructions,
  online logbooks, instructions and information created by participants
  will be merged into a comprehensive website that will be available to
  the community and the public for self-guided training or for use in
  various educational settings (for example, a graduate course at a
  university could assign some of the projects as homework or an extra
  credit project, etc).
\end{itemize}


\subsection{Course content, learning outcomes and detailed plan}

The course will be taught as an intensive course of duration of three
weeks, with a total time of 45 h of lectures, 45 h of exercises and a
final assignment of 2 weeks of work. The total load will be
approximately 160-170 hours, corresponding to \textbf{7 ECTS} in
Europe.  The final assignment will be graded with marks A, B, C, D, E
and failed for Master students and passed/not passed for PhD
students. A course certificate will be issued for students requiring
it from the University of Trento.

The organization of a typical course day is as follows:

\begin{enumerate}
\def\labelenumi{\arabic{enumi}.}
\item
  9am-12pm: Lectures, project relevant information and directed
  exercises
\item
  12pm-2pm: Lunch
\item
  2pm-6pm: Computational projects, exercises and hands-on sessions
\item
  6pm-7pm: Wrap-up of the day and eventual student presentations
\end{enumerate}

If approved by the ECT* board of directors, our preferred time slot
would be from the second half of June till the second half of July.


\subsubsection{First week, schedule and learning outcomes}

\begin{table}[hbtp]
\begin{tabular}{|l|l|l|l|l|} \hline
    & First session  & Second session  & Exercises and project work & Student presentations \\ \hline
  Monday & Review of Linear Algebra and & Qubits, measurements and  & Work on codes for basis & \\
         & density matrices and states & quantum gates and circuits & and quantum gates & \\
  Tuesday & Quantum Fourier & Quantum phase  & QFT exercises &  \\
          & transforms (QFT) & estimation algorithm (QPE) & and codes & \\
  Wednesday  & Quantum algorithms &  Quantum advantage & Work on exercises & \\    
  Thursday & Variational Quantum & Simple Hamiltonians & Implementing QPE & \\
            & Eigensolver (VQE) & & Work  on VQE & \\
  Friday  & Nuclear physics  & Lipkin model & VQE implementation & \\
           & Hamiltonians     & and Rodeo algorithm  & of the Lipkin model & \\ \hline
\end{tabular}
\caption{Teaching schedule first week}
\end{table}

The first week focuses, after a reminder of central linear algebra
elements, on basic ingredients of quantum computing such as rewriting
quantum mechanical operations as quantum gates and circuits, how to
perform measurements and how to obtain eigenvalues of selected
Hamiltonians. We will also (Wednesday) discuss some selected quantum
algorithms, such as Grover's algorithm, Simon's algorithms and quantum
advantage via Shor's algorithm.


To obtain the eigenvalues we will discuss the quantum
phase estimation algorithm (which requires a discussion of Quantum
Fourier transforms) and the widely used variational quantum
eigensolver (VQE). After having introduced some simpler Hamiltonians
defined by various Pauli matrices, we will demonstrate how to rewrite
a widely used Hamiltonian given by a second-quantized representation
in terms of various Pauli matrices. The Hamiltonian we will focus on
the first week is Lipkin Hamiltonian which does not require
a so-called Jordan-Wigner transformation. This transformation will
be discussed during the second week.  The students will work on
analytical exercises as well as computational exercises. The latter
will focus on developing a code which implements the VQE method for
finding the eigenvalues of the above Hamiltonians.  The Rodeo algorithm will also be discussed. This code can be
extended upon and can be used to define a final project students can
hand in for final credits.

Many of the topics discussed during the first week, will serve as
background material for the next two weeks.

The schedule for the student presentations will be finalized during
the Talent course.

\subsubsection{Second week, schedule and learning outcomes}


\begin{table}[hbtp]
\begin{tabular}{|l|l|l|l|l|} \hline
& First session  & Second session  & Exercises and project work & Student presentations \\ \hline
  Monday  & Hamiltonian dynamics & Lie-Trotter-Suzuki & Work on VQE& \\
          &                      & formulas for simulations  & and exercises & \\ 
  Tuesday & Encoding (relativistic and nonrelativistic)  & Encoding & Exercises & \\
          & fermions and bosons on quantum computers     & second part & and project work & \\              
  Wednesday & Quantum simulation & Quantum simulation  & Simulations  & \\
            & of pionless EFT    & of pionless EFT, part 2 & of pionless EFT & \\
  Thursday & Nuclear response & Neutrino dynamics  & Exercises and & \\
           & functions  &  in dense environments & project work & \\
  Friday & Noise mitigation and & Quantum error correction   & Exercises and project work & \\ 
  & NISQ computing & fault tolerance  &  & \\ \hline  
\end{tabular}
\caption{Tentative schedule second  week}
\end{table}

The second week starts with a discussion of product formulae such as
the Lie-Trotter-Suzuki approximation and how to simulate Hamiltonian
dynamics.  Thereafter we discuss in detail how to encode fermionic and
bosonic systems through for example the so-called Jordan-Wigner
transformation. This will allow us to study more general Hamiltonians
such as a pionless EFT based Hamiltonian and nuclear response function
and neutrino dynamics.  Entanglement in nuclear many-body systems will
also be discussed before we wrap up the week with a discussion of
noise mitigation and quantum error correction algorithms.

Given the more general Hamiltonians, the codes developed during the
first week can be extended to include more Hamiltonians and systems of
relevance for nuclear physics.


\subsubsection{Third week, schedule and learning outcomes}


\begin{table}[hbtp]
\begin{tabular}{|l|l|l|l|l|} \hline
    & First session  & Second session  & Exercises and project work & Student presentations \\ \hline
  Monday & Quantum simulation of   & Quantum simulation of   & Exercises & \\
         & scattering in scalar field theory, I & scattering in scalar field theory, II & and project work & \\
  Tuesday & Hamiltonian formulations of & Hamiltonian formulations of & Exercises & \\
          & gauge theories (Abelian)    & gauge theories (non-Abelian) & and project work & \\
  Wednesday & Time evolution  & Time evolution & Exercises  & \\
            & in gauge theories & in gauge theories & and project work \\
  Thursday & Applications to QFT & Applications to QFT& Project work & \\
  Friday & Summary of course & Discussion of projects  & Project work & \\ \hline
\end{tabular}
\caption{Teaching schedule third  week}
\end{table}

Depending on our progress during the first two weeks, the schedule of
the last week may be subject to changes. The tentative plan for the
final week is to dedicate it a discussion of quantum computing for
quantum (gauge) field theories. Here the learning outcomes will focus
on quantum simulations of scattering in scalar field theory,
Hamiltonian formulations of gauge theories and time evolution in gauge
theories, Abelian and non-Abelian, with final applications. The course
ends with a summary and discussions of the projects.




\subsection{Instructors and organizers}\label{instructors-and-organizers}

The organizers and instructors are


\begin{enumerate}
\item 
\href{https://directory.natsci.msu.edu/Directory/Profiles/Person/101033}{Alexei
Bazavov} at Department of Computational Mathematics, Science and
Engineering and Department of Physics and Astronomy, Michigan State
University, East Lansing, MI 48824, USA. He is a theoretical particle physicist specializing in study of strongly coupled theories, in particular, Quantum Chromodynamics. Items of particular interest to him include quantum computing and algorithms, quantum field theory, finite-temperature field theory, lattice gauge theory with applications to particle and nuclear physics, parallel algorithms, iterative solvers, molecular dynamics algorithms, inverse problems and Bayesian inference, ultra-cold atomic systems and quantum simulations and effective field theory.

\item
  \href{https://umdphysics.umd.edu/people/faculty/current/item/927-davoudi.html}{Zohreh
  Davoudi}, Department of Physics and Center for Quantum Information and Computer Science (QuICS), University of Maryland, College Park,
  MD 20742, USA. Davoudi is an expert in lattice QCD for nuclear physics. She further
  specializes in quantum simulation and computing for quantum field
  theories, including lattice gauge theories and effective field
  theories of nuclear physics.

\item
  \href{http://mhjgit.github.io/info/doc/web/}{Morten Hjorth-Jensen} at
  Facility for Rare Isotope Beams and Department of Physics and
  Astronomy, Michigan State University, East Lansing, MI 48824, USA \&
  Department of Physics, University of Oslo, N-0316 Oslo, Norway. Hjorth-Jensen has his background in studies of different many-body theories applied to problems in nuclear physics and condensed matter physics. He works also on quantum engineering and machine learning applied to many-body systems. He has over many years developed introductory and advanced learning material in many-body physics, quantum computing, computational physics and machine learning.
\item
  \href{https://frib.msu.edu/for-students/faculty/lee-profile}{Dean Lee}
  at Facility for Rare Isotope Beams and Department of Physics and
  Astronomy, Michigan State University, East Lansing, MI 48824, USA. His research is focused on connecting fundamental physics to forefront experiments. He studies many aspects of quantum few- and many-body systems. Together with collaborators, he has developed lattice Monte Carlo methods that probe strongly-interacting systems and study superfluidity, nuclear clustering, phase transitions, and other emergent phenomena from first principles. He  is also engaged in novel applications of new technologies for scientific research. This includes new algorithms for quantum computing and the development of emulators and machine learning algorithms based on concepts such as eigenvector continuation.
\item
  \href{https://www.ryanlarose.com/}{Ryan LaRose} at at Department of
  Computational Mathematics, Science and Engineering, Michigan State
  University, East Lansing, MI 48824, USA. He does research in computational physics and quantum information science. He is  interested in both the physics of computation and the computation of physics - that is, what quantum physics can tell us about information and computer science, and how quantum computers can solve practical problems in physics and related fields.
\item
  \href{https://webapps.unitn.it/du/en/Persona/PER0016084/Didattica}{Alessandro
  Roggero}, Department of Physics, University of Trento, Povo, 38123
  Trento, Italy. He works on  simulations of strongly correlated quantum many-body systems using a combination of classical techniques and calculations on
quantum devices. In particular, in understanding the role of entanglement in both observable properties of physical systems and as a fundamental tool to understand the structure of many-body states. His research interests include quantum simulations of inelastic nuclear processes, neutrinos in dense matter, quantum computing and algorithms, quantum Monte Carlo methods and effective field theory.

\item
  An eventual Post-doctoral fellow and an advanced graduate student as
  teaching assistants.
\end{enumerate}

Morten Hjorth-Jensen and Alessandro Roggero will also function as student
advisors and coordinators.

\subsection{Audience and Prerequisites}
The potential participants are
Students and post-doctoral fellows interested in quantum computing
applied to nuclear physics problems, from nuclear structure problems to quantum
field theories. The material will be of interest, and accessible, to
both theorists and experimentalists, and will include learning the
practical use of quantum computing software in order to interpret and
study nuclear systems.

The students are expected to have operating programming skills in
compiled programming languages like Fortran or C++ or preferentially in
an interpreted language like Python and knowledge of quantum mechanics
at an intermediate level.

\subsection{Admission}

The target group is Master of Science students, PhD students and early
post-doctoral fellows. Also senior staff can attend but they have to be
self-supported. The maximum number of students is 20-30, of which
hopefully 15-20 can receive full local support.

The process of selections of the students will be managed in agreement
with the ECT*.

\subsection{Preliminary budget}

We expect to accept between 20-30 students. Local students from the
University of Trento are fully self-supported. If approved, we would
very much appreciate if the ECT* can sponsor 15-20 of the selected
students with local expenses, that is lodging and meals during weekdays.
Any additional funds for sponsoring further students is highly
appreciated.

All travel expenses will be covered by the respective home institute.
Instructors are self-supported. We plan to raise additional funds to
cover local support for additional students and the expenses of the
instructors.

There is no participation fee. Administrative support from the ECT* in
organizing the course and setting up the application procedure is
essential for a smooth (as always) outcome. The administrative
experience of the staff at the ECT* has been unique and essential in
running successfully our previous Talent courses. We would thus highly
appreciate it if these services are provided if the proposal is
approved.



\end{document}










